\par\Needspace{6\baselineskip}\medskip
\begin{center}\hypertarget{native-192-article143}{}\pdfbookmark[2]{143.}{article-143}143.\end{center}
As soon as the angle $z$ has been found, $r'$ is given at once by the equation
\[
r'=\frac{R'\sin\delta'}{\sin z}\mathpunct{\cdot}
\]
Further, from equations $P=\dfrac{n''}{n}$ and III we obtain
\[
\frac{n'r'}{n}=\frac{(P+a)R'\sin\delta'}{b\sin(z-\sigma)}
\]
\[
\frac{n'r'}{n''}=\frac{1}{P}\cdot\frac{n'r'}{n}\mathpunct{\cdot}
\]

\SourceLeaf{193}We now treat the formulas for determining the positions of the points $C,C''$ from the position of $C'$ in a way that makes their general validity immediately clear, even for cases that Fig. 4 does not show. We observe that the sine of the distance of $C'$ from the great circle $CB$ (the distance being taken as positive in the upper region and negative in the lower) equals the product of $\sin\epsilon''$ and the sine of the distance of $C'$ from $D''$ measured in the forward direction, and is therefore $=-\sin\epsilon''\sin C'D''=-\sin\epsilon''\sin(z+A'D''-\delta')$. Likewise, the sine of the distance of $C''$ from the same great circle is $=-\sin\epsilon'\sin C''D'$. These sines are plainly as $\sin CC'$ to $\sin CC''$, or as $\dfrac{n''}{rr'}$ to $\dfrac{n'}{rr''}\mathpunct{\raisebox{.45ex}{,}}$ or as $n''r''$ to $n'r'$. Thus, on setting $C''D'=\zeta''$, we have
\[
\tag*{V.}r''\sin\zeta''=\frac{n'r'}{n''}\cdot\frac{\sin\epsilon''}{\sin\epsilon'}\sin(z+A'D''-\delta').
\]
In exactly the same way, on setting $CD'=\zeta$, we obtain
\[
\tag*{VI.}r\sin\zeta=\frac{n'r'}{n}\cdot\frac{\sin\epsilon}{\sin\epsilon'}\sin(z+A'D-\delta')
\]
\[
\tag*{VII.}r\sin(\zeta+AD''-AD')=r''P\cdot\frac{\sin\epsilon}{\sin\epsilon''}\sin(\zeta''+A''D-A''D').
\]
By combining equations V and VI with the following equations taken from Article 139,
\[
\tag*{VIII.}r''\sin(\zeta''-A''D'+\delta'')=R''\sin\delta''
\]
\[
\tag*{IX.}r\sin(\zeta-AD'+\delta)=R\sin\delta,
\]
the quantities $\zeta,\zeta'',r,r''$ can be derived according to the rule in Article 78. To make this calculation easier, we give the formulas themselves here. Set
\[
\tag*{17)}\frac{R\sin\delta}{\sin(AD'-\delta)}=\varkappa
\]
\[
\tag*{18)}\frac{R''\sin\delta''}{\sin(A''D'-\delta'')}=\varkappa''
\]
\[
\tag*{19)}\frac{\cos(AD'-\delta)}{R\sin\delta}=\lambda
\]
\[
\tag*{20)}\frac{\cos(A''D'-\delta'')}{R''\sin\delta''}=\lambda''.
\]
The calculation of these quantities, or rather of their logarithms, is still independent of $P$ and $Q$ and is to be regarded as the \emph{fifth} and last of what may be called the preliminary operations. It is conveniently carried out at once with the calculation of $a,b$, that is, with the fourth task, where $a=\dfrac{\varkappa}{\varkappa''}\mathpunct{\cdot}$ — If we then \SourceLeaf{194}set
\[
\begin{aligned}
\frac{n'r'}{n}\cdot\frac{\sin\epsilon}{\sin\epsilon'}\cdot\sin(z+A'D-\delta')&=p\\[.5ex]
\frac{n'r'}{n''}\cdot\frac{\sin\epsilon''}{\sin\epsilon'}\cdot\sin(z+A'D''-\delta')&=p''\\[.5ex]
\varkappa(\lambda p-1)&=q\\[.5ex]
\varkappa''(\lambda''p''-1)&=q''
\end{aligned}
\]
we obtain $\zeta$ and $r$ from $r\sin\zeta=p$, $r\cos\zeta=q$, and $\zeta''$ and $r''$ from $r''\sin\zeta''=p''$, $r''\cos\zeta''=q''$. There can be no ambiguity here in determining $\zeta$ and $\zeta''$, since $r$ and $r''$ must necessarily turn out to be positive. The completed calculation can be checked by equation VII if desired.

There are, however, two cases in which a different method must be followed. Whenever the point $D'$ coincides with $B$ or is opposite it on the sphere, that is, whenever $AD'-\delta=0$ or $=180^0$, equations VI and IX must necessarily be identical, giving $\varkappa=\infty$, $\lambda p-1=0$, and hence an indeterminate $q$. In this case $\zeta''$ and $r''$ are determined as described above, and then $\zeta$ and $r$ must be found by combining equation VII with VI or IX. We forgo reproducing the formulas themselves, which can be taken from Article 78. We merely note that the same method is preferable even when $AD'-\delta$ is neither $=0$ nor $=180^0$ but is a very small arc, because the earlier method would then lack the required precision. In particular, we combine equation VII with VI or IX according as $\sin(AD''-AD')$ is greater or less than $\sin(AD''-\delta)$.

Likewise, when the point $D'$ or its opposite coincides with $B''$ or is only a little way from it, determining $\zeta'',r''$ by the preceding method would be either impossible or unreliable. In that case $\zeta$ and $r$ are determined by that method, and then $\zeta''$ and $r''$ are found by combining equation VII with V or VIII, according as $\sin(A''D-A''D')$ is greater or less than $\sin(A''D-\delta'')$. There is, however, no need to fear \emph{simultaneous} coincidence or near coincidence of $D'$ with the points $B,B''$ or their opposites. We have already excluded the case in which $B$ coincides with $B''$ or is very close to it from our discussion above in Article 138.
